% ============================================
% CONT-02 (ZH)
% 中文翻译 — SPEC-935-R356
% ============================================
\section*{\textbf{CONT-02} --- 您当时就在那里}
\markboth{CONT-02 --- 您当时就在那里}{}

\noindent \textit{您一直都在}

\noindent 您以为自己在阅读关于过去的事。

\noindent 您以为若阿金是一个人物。乱葬坑是一个场景。黄眼是创伤的隐喻。那生物是他心灵的投射、一种心理防御、一种解离机制。

\noindent 您错了。

\noindent 您当时就在那座乱葬坑里。

\noindent 不是物理意义上。您不是1907年出生的，没有走过基谢拉莫宾的塞尔唐，没有亲眼看着父亲饿死。这一切都是真的。

\noindent 但您以另一种方式在那里。

\noindent 您在那里，是因为若阿金的故事，就是每一个曾经感受过饥饿之人的故事。不是食物的饥饿。是对这个世界无法给予之物的饥饿。对意义的饥饿。对情感的饥饿。对正义的饥饿。对被看见的饥饿。

\noindent 您感受过这种饥饿。

\noindent 您曾在人群中感到孤独。曾觉得内心有什么不对劲。某种没有名字的东西。某种您用消遣、用工作、用噪音、用任何能阻止寂静的东西喂养的东西。曾在半夜醒来，觉得房间里并非只有您一人。

\noindent 感受寒意从后颈一直爬升到头皮。您的胃在您想起原因之前，先收紧了一下。

\noindent 若阿金在乱葬坑里遇到的生物，并不是超自然实体。

\noindent 它是每个人都背负的饥饿。区别在于，若阿金\textit{看见了}它。触碰了它。让它进来了。

\noindent 它也已经在您体内了。

\noindent 您只是还没有看见。

\noindent 但在读完这本书之后，您会开始看见。

\noindent 在房间的角落里。

\noindent 在玻璃的倒影中。

\noindent 在街上那条狗的黄色眼睛里。

\noindent 而某个夜晚，当您躺下、灯已熄灭时，您会感到后颈一阵沉重。一股温热的气息。

\noindent 别动。

\noindent 您不会害怕。您会认出它。

\noindent 您会想：\textit{\textquotedblleft{}您当时也在那里。\textquotedblright{}}

\noindent 而那生物会回答：

\noindent \textit{\textquotedblleft{}我从未离开。\textquotedblright{}}

\vspace{0.5em}
\noindent\textit{\textbf{↪ 路线:} \rota{CONT-03} \rotasep \rota{MEM-04}}
